\question{2.m5}{Voor een steekproef van 500 werknemers van een groot concern is in 2001 de verdeling van de jaarinkomens onderzocht.
Hierbij bleek dat de standaarddeviatie van de inkomens 7.500 guldens bedroeg.
Ter wille van een internationale publicatie moeten de gegevens worden uitgedrukt in euro's
(1 euro = afgerond 2.20 gulden). De standaarddeviatie gemeten in euro's bedraagt dus ...}
\begin{enumerate}[label=(\alph*)]
    \item 16.500
    \item 36.300
    \item 3.409
    \item 1.100
\end{enumerate}
\answer{
    De standaarddeviatie in euro's wordt berekend door de standaarddeviatie in guldens te delen door de wisselkoers:
    \[
        s_{\textrm{euro}} = \frac{s_{\textrm{gulden}}}{2,20} = \frac{7.500}{2.20} \approx 3.409.
    \]
    Het juiste antwoord is dus (c). 
}